
\documentclass[../../main.tex]{subfiles}
\begin{document}

% 年度の大きな表示
\begin{center}
  {\fontsize{48pt}{58pt}\selectfont \textbf{1980}\normalsize\textbf{年度}}
\end{center}

\vspace{10mm}

% 本文

\vspace{15mm}

% 出題テーマと難易度の表
\noindent\textbf{出題テーマと難易度}

\vspace{5mm}

\begin{center}
  \shadowbox{%
    \begin{tabular}{|c|p{8cm}|c|}
      \hline
      \textbf{問題番号} & \textbf{テーマ} & \textbf{難易度} \\
      \hline
      第1問 & 行列と1次変換 & \\
      \hline
      第2問 & 図形と角の二等分線 & \\
      \hline
      第3問 & 指数関数と接線の本数 & \\
      \hline
      第4問 & 定積分と数列の証明 & \\
      \hline
    \end{tabular}%
  }
\end{center}

\vspace{15mm}

% 解答
\noindent\textbf{解答}

\vspace{5mm}

\begin{center}
  \shadowbox{%
    \renewcommand{\arraystretch}{1.6}
    \begin{tabular}{|c|c|p{8cm}|}
      \hline
      \textbf{問題番号} & \textbf{小問} & \textbf{答え} \\
      \hline
      第1問 & - &  \\
      \hline
      第2問 & - & $\dfrac{2}{3}<x<\sqrt{2}$ \\
      \hline
      第3問 & - & $b\le0$または$b=e^a$のとき$1$個／$0<b<e^a$のとき$2$個／$b>e^a$のとき$0$個 \\
      \hline
      \multirow{3}{*}{第4問} & (1) & 証明問題 \\
      \cline{2-3}
                             & (2) & 証明問題 \\
      \cline{2-3}
                             & (3) & 証明問題 \\
      \hline
    \end{tabular}%
    \renewcommand{\arraystretch}{1.0}
  }
\end{center}

\vspace{15mm}

% 解答時間
\noindent\textbf{試験形態}

\vspace{5mm}

\begin{center}
  \shadowbox{%
    \begin{tabular}{p{8cm}}
      （要確認）
    \end{tabular}%
  }
\end{center}

\end{document}
